\par
\underline{证明}。设 $p:S'\longrightarrow S$ 是 $S$ 的双有理变换。将 $D$ 与 $E$ 的全变换记为 $D^*$ 与 $E^*$，将其严格变换记为 $D'$ 与 $E'$。

%%%%%%%%%% source scan idx 25 | folio 18 %%%%%%%%%%

\underline{引理} 2.9.1. (i) \underline{全变换 $D^*$ 由下列条件刻画：}

(a)\quad \underline{$D^*-D'$ 的支集包含于 $p^{-1}(s)$；}

(b)\quad \underline{若 $X$ 的支集包含于 $p^{-1}(s)$，则 $(D^*,X)=0$。}

(ii)\quad $(D,E)=(D^*,E^*)=(D^*,E')=(D',E^*)$.

断言 (i)(a) 显然成立，(i)(b) 则由 1.6(vi) 得到。由 1.9 可知 (a) 与 (b) 确定 $D^*$。由构造，映射 $\underline{O}_E\longrightarrow p_*\underline{O}_{E'}$ 的余核为有限长度模；所以

\begin{equation*}
\begin{aligned}
(D^*,E')
&=\chi(Lp^*\underline{O}_D\otimes\underline{O}_{E'})\\
&=\chi(\underline{O}_D\otimes Rp_*\underline{O}_{E'})
=\chi(\underline{O}_D\otimes\underline{O}_E)
=(D,E),
\end{aligned}
\end{equation*}
结合 (i) 便推出 (ii)。

取 $S'$ 为 $S$ 在 $s$ 处的爆破。由（2.9.1）、（2.6.3）和 2.7，

\begin{equation}\tag{2.9.2}
\begin{aligned}
(D,E)
&=(D^*,E')
=(D'+\mu_s(D)s^*,E')\\
&=(D',E')+\mu_s(D)(s^*,E')\\
&=(D',E')+\mu_s(D)\mu_s(E)[k(s):k].
\end{aligned}
\end{equation}
现在对 $(D,E)$ 归纳，反复应用（2.9.2），即得定理 2.9。

(2.10) 设 $A$ 是一维约化 Noether 局部环，$s$ 是 $\operatorname{Spec}(A)$ 的闭点，$\widetilde A$ 是 $A$ 的正规化，$k$ 是使 $k(s)$ 在其上有限的域。假设 $\widetilde A$ 在 $A$ 上有限（等价地，完备化 $\widehat A$ 约化）。$\operatorname{Spec}(A)$ 在 $s$ 处的\underline{奇点的 $\delta$-不变量}，定义为 $A$-模 $\widetilde A/A$ 在 $k$ 上的长度。

%%%%%%%%%% source scan idx 26 | folio 19 %%%%%%%%%%

\underline{定理} 2.11. \underline{设 $D$ 是 $S$ 上的约化曲线。则：}

(i)\quad \underline{对任意 $t\in S^*$，沿用 2.8 的记号，有}

\begin{equation*}
\mu_t(D)
=\sum_{\substack{t'\ \text{邻接}\\\text{于}\ t}}
\mu_{t'}(D)[k(t'):k(t)] \quad .
\end{equation*}

(ii)\quad \underline{假设 $D$ 的正规化 $\widetilde D$ 在 $D$ 上有限，令 $\delta=\chi(\underline{O}_{\widetilde D}/\underline{O}_D)$ 为 $D$ 在 $s$ 处的奇点的不变量。除有限多个 $t\in S^*$ 外，$\mu_t(D)$ 均为零或一，且}

\begin{equation*}
\delta
=\sum_{t\in S^*}\frac12\mu_t(D)\bigl(\mu_t(D)-1\bigr)[k(t):k] \quad .
\end{equation*}

设 $p:S_1\longrightarrow S$ 是 $S$ 沿 $\{s\}$ 的爆破，$D^*$ 是 $D$ 的全变换，$D'$ 是其严格变换。由（2.7），

\begin{equation}\tag{2.11.1}
D^*=D'+\mu_s(D)s^*
\end{equation}
且
\begin{equation}\tag{2.11.2}
\mu_s(D)=(s^*,D')\,[k(s):k]^{-1} \quad .
\end{equation}
小曲线 $s^*$ 无奇点。因此，由 2.8 和定义 2.5，将公式 2.9 应用于 $S_1$ 上的除子 $s^*$ 与 $D'$，便成为

\begin{equation}\tag{2.11.3}
(s^*,D')
=\sum_{\substack{t'\ \text{邻接}\\\text{于}\ s}}
\mu_{t'}(D)[k(t'):k] \quad .
\end{equation}
当 $s=t$ 时，断言 (i) 由（2.11.2）和（2.11.3）得到；将此结果应用于 $S(t)$ 在 $t$ 处的局部化（2.3），即得一般情形。

设 $m$ 是定义 $\{s\}$ 的极大理想，$n\geq0$。已知（SGA VII 3.5）

\begin{equation}\tag{2.11.4}
p_*\underline{O}_{S_1}(-ns^*)=m^n
\quad\text{且}\quad
R^ip_*\underline{O}_{S_1}(-ns^*)=0\quad(i>0).
\end{equation}

%%%%%%%%%% source scan idx 27 | folio 20 %%%%%%%%%%

由此以及正合列

\begin{equation*}
0\longrightarrow\underline{O}_{S_1}(-ns^*)
\longrightarrow\underline{O}_{S_1}
\longrightarrow\underline{O}_{ns^*}
\longrightarrow0
\end{equation*}
得到

\begin{equation}\tag{2.11.5}
p_*\underline{O}_{ns^*}=\underline{O}_S/m^n
\quad\text{且}\quad
R^ip_*\underline{O}_{ns^*}=0\quad(i>0),
\end{equation}
以及

\begin{equation}\tag{2.11.6}
\chi(\underline{O}_{ns^*})
=\chi(\underline{O}_S/m^n)
=\frac{n(n+1)}2[k(s):k] \quad .
\end{equation}
在（2.11.4）中取 $n=0$。对 $S$ 上的任一凝聚层复形 $K$，投影公式给出

\begin{equation*}
K\xrightarrow{\;\sim\;}Rp_*Lp^*K
\qquad
\bigl(\simeq K\overset{\mathbb L}{\otimes}Rp_*\underline{O}_{S_1}\bigr).
\end{equation*}
特别地，由于 $Lp^*\underline{O}_D=p^*\underline{O}_D=\underline{O}_{D^*}$，有

\begin{equation}\tag{2.11.7}
p_*\underline{O}_{D^*}=\underline{O}_D
\quad\text{且}\quad
R^ip_*\underline{O}_{D^*}=0\quad(i>0).
\end{equation}
置 $E=\mu_s(D)s^*$。由（2.11.1）、（2.9.1）(i)(b) 和（2.6.1），

\begin{equation}\tag{2.11.8}
\begin{aligned}
\chi(\underline{O}_{D'\cap E})
&=(D',E)=(D^*-E,E)=-(E,E)\\
&=-\mu_s(D)^2(s^*,s^*)
=\mu_s(D)^2[k(s):k].
\end{aligned}
\end{equation}
正合列

\begin{equation*}
0\longrightarrow\underline{O}_{D^*}
\longrightarrow\underline{O}_{D'}\oplus\underline{O}_E
\longrightarrow\underline{O}_{D'\cap E}
\longrightarrow0
\end{equation*}
取正向像并利用（2.11.7）后，给出正合列

\begin{equation*}
0\longrightarrow\underline{O}_D
\longrightarrow p_*\underline{O}_{D'}\oplus p_*\underline{O}_E
\longrightarrow p_*\underline{O}_{D'\cap E}
\longrightarrow0,
\end{equation*}
从而

\begin{equation*}
0\longrightarrow p_*\underline{O}_E
\longrightarrow p_*\underline{O}_{D'\cap E}
\longrightarrow p_*\underline{O}_{D'}/\underline{O}_D
\longrightarrow0 \quad .
\end{equation*}

%%%%%%%%%% source scan idx 28 | folio 21 %%%%%%%%%%

所以，由（2.11.6）和（2.11.8），

\begin{equation}\tag{2.11.9}
\begin{aligned}
\chi(p_*\underline{O}_{D'}/\underline{O}_D)
&=\mu_s(D)^2[k(s):k]
-\frac12\mu_s(D)\bigl(\mu_s(D)+1\bigr)[k(s):k]\\
&=\frac12\mu_s(D)\bigl(\mu_s(D)-1\bigr)[k(s):k] \quad .
\end{aligned}
\end{equation}
概形 $D$ 与 $D'$ 均约化，只在 $s$ 上方不同，且 $D'$ 在 $D$ 上有限。故它们具有相同的正规化，$D$ 在 $s$ 处的奇点的 $\delta$-不变量，是 $\chi(p_*\underline{O}_{D'}/\underline{O}_D)$ 与 $D'$ 在 $D'\cap s^*$ 各点处的奇点不变量 $\delta$ 之和。对 $\delta$ 归纳证明 2.11(ii)。若 $\mu_s(D)>1$，则归纳假设可应用于 $D'$，由（2.11.9）即得 2.11(ii)。若 $\mu_s(D)=1$，则 $(D',s^*)=1$，且 $D'\cap s^*$ 是 $D'$ 上的正则除子；由 EGA IV 19.1.1，$D'$ 正则。公式（2.11.9）于是表明 $D$ 本身也正则，再由 2.8 即得 2.11(ii)。

